\section{Conditional theory for interface repair}
\label{app:composite-theory}

This appendix separates algebraic identifiability, the authority of the chosen
action interface, and stability of the implemented observer. Its dynamical
results concern an explicitly assumed local model. They do not turn the fitted
healthy reference or the reported per-step eigenvalue screens into a certificate
for the physical VLA--robot loop.

\subsection{Identifying a disturbance is different from cancelling its effect}

Let a constant parameter $f\in\mathbb R^p$ produce a bias-corrected observation
$z=Hf$ and an unwanted physical response $Gf\in\mathbb R^n$. A correction
$c\in\mathbb R^k$ removes the response $Dc$, where $D\in\mathbb R^{n\times k}$
describes the available action interface. Thus $H$ describes observation and $D$
describes correction authority; neither is assumed equal to $G$.

\begin{proposition}[Identifiability and exact linear cancellation]
\label{prop:authority}
The parameter $f$ is uniquely identifiable from $z$ if and only if
$\ker H=\{0\}$. There exists a linear rule $c=Fz$ cancelling every parameter,
$DFH=G$, if and only if
\begin{equation}\label{eq:authority-factorization}
\ker H\subseteq\ker G,\qquad
\operatorname{range}G\subseteq\operatorname{range}D.
\end{equation}
With corrections restricted to a set $\mathcal U$, exact cancellation for a
particular known $f$ is feasible if and only if $Gf\in D\mathcal U$.
\end{proposition}
\begin{proof}
Two parameters yield the same observation exactly when their difference lies
in $\ker H$, proving the first assertion. If $DFH=G$, then $Hv=0$ implies
$Gv=0$, and every column of $G$ lies in the range of $D$, proving necessity.
Conversely, let $H^\dagger,D^\dagger$ denote Moore--Penrose inverses. The two
conditions give $GH^\dagger H=G$ and $DD^\dagger G=G$, respectively. Therefore
$F=D^\dagger GH^\dagger$ satisfies $DFH=G$. The bounded assertion is precisely
the existence of $c\in\mathcal U$ with $Dc=Gf$.
\end{proof}

Full parameter identifiability is consequently neither sufficient nor necessary
for cancellation of the physical effect. For example, $H=I_2$, $G=I_2$ and
$D=(1,0)^\top$ make both fault coordinates observable while leaving the second
uncorrectable. Conversely, $H=(1,0)$, $G=(1,0)$ and $D=1$ permit exact cancellation
although the irrelevant second fault coordinate is unidentifiable. These are
local interface statements: configuration dependence alone establishes neither
of the range conditions. Unknown output bias further limits identifiability as
in Proposition~\ref{prop:sep}.

\subsection{The recursion implemented by the composite observer}

For the following result the estimate and fault have the same $p$ coordinates.
Suppose that, over a stated operating region, actual-minus-reference position
error obeys
\begin{equation}\label{eq:local-position-error}
e_{t+1}=Ae_t+Gf_t-D\hat\theta_t+w_{t+1},\qquad
z_{t+1}=Hf_t+v_{t+1}.
\end{equation}
Here $w$ includes model error and $v$ includes observation error, remaining bias,
and any discrepancy in the calibrated residual relation. Cancellation by the
estimate $\hat\theta=f$ requires $(G-D)f=0$, for example when $G=D$ or the
admissible faults occupy only coordinates where the two maps agree. The more
general authority condition in Proposition~\ref{prop:authority} can instead
require a different allocation of the identified parameter. For ALOHA,
$H=M$ is the full residual sensitivity,
$D=B\operatorname{diag}(m)$ includes the correction mask, and the fitted reference
is driven by the live \emph{raw} action. Subtracting this reference cancels the
common raw input in the assumed linear model; it does not construct the healthy
counterfactual trajectory of the nonlinear policy.
Using the fitted $D$ in the physical recursion is an assumption: if the actual
map is $D_{\rm phys}$, the remainder must include
$-(D_{\rm phys}-D)\hat\theta_t$. Applying the bounds below then requires a
bound on this remainder over the operating region; fitting $D$ does not supply
that bound automatically.

Write $a=\exp(-\lambda\Delta t)$ and
$h=\Delta t$ when $\lambda=0$, otherwise $h=(1-a^2)/(2\lambda)$.
The implementation uses
\begin{align}
P_t^-&=a^2P_t+hQ_c,
&K_t&=P_t^-H^\top(HP_t^-H^\top+R)^{-1},\\
P_t^+&=(I-K_tH)P_t^-(I-K_tH)^\top+K_tRK_t^\top,
&P_{t+1}&=P_t^+,\\
T_t&=\Delta t\,\eta P_t^+D^\top L,
&\bar\theta_{t+1}&=a\hat\theta_t+K_t(z_{t+1}-Ha\hat\theta_t)+T_te_{t+1},\\
\hat\theta_{t+1}&=\Pi_{\mathcal U}(\bar\theta_{t+1}).&&
\label{eq:appendix-composite-recursion}
\end{align}
Here $Q_c$ is process covariance per unit time, $R\succ0$ is covariance per
observation, and $L\succ0$ is the chosen position metric. The action has already
used $\hat\theta_t$ before $e_{t+1}$ and $z_{t+1}$ are measured. $P_t^+$ is the
posterior covariance of the specified prediction-only statistical model; after
adding tracking feedback it is a design preconditioner, not automatically the
true covariance of the composite estimation error.

Let $\delta_t=\hat\theta_t-f_t$, $C_t=I-K_tH$, and
$\xi_t=(e_t^\top,\delta_t^\top)^\top$. Define
\begin{align}
b_{t+1}&=(G-D)f_t+w_{t+1},
&\Delta f_t&=f_{t+1}-f_t,\\
p_{t+1}&=\Pi_{\mathcal U}(\bar\theta_{t+1})-\bar\theta_{t+1}.
\end{align}
Direct substitution yields the exact error recursion under
\eqref{eq:local-position-error}:
\begin{equation}\label{eq:composite-augmented}
\begin{aligned}
\xi_{t+1}&=\mathcal A_t\xi_t+d_{t+1},\\
\mathcal A_t&=
\begin{bmatrix}
A&-D\\
T_tA&aC_t-T_tD
\end{bmatrix},\\
d_{t+1}&=
\begin{bmatrix}
b_{t+1}\\
T_tb_{t+1}+(a-1)C_tf_t-\Delta f_t+K_tv_{t+1}+p_{t+1}
\end{bmatrix}.
\end{aligned}
\end{equation}
In particular, the lower-left block is $T_tA$, not $T_t$: tracking uses the
\emph{new} position error. The positive tracking sign is also deliberate. For
fixed $e_t,f_t$, the gradient of
$\frac12\|Ae_t+Gf_t-D\theta\|_L^2$ at $\theta=\hat\theta_t$, in the noiseless
model, is $-D^\top Le_{t+1}$. Hence tracking alone is a preconditioned negative
gradient step. This direction does not establish stability of the full finite-step
observer--controller recursion.

\begin{proposition}[Conditional contraction and disturbance bound]
\label{prop:composite-iss}
Assume \eqref{eq:local-position-error} holds along the trajectory. Suppose a
\emph{single} matrix $V\succ0$ and $0<\rho<1$ satisfy
\begin{equation}\label{eq:common-composite-metric}
\mathcal A_t^\top V\mathcal A_t\preceq\rho^2 V
\quad\text{for every gain encountered.}
\end{equation}
Then, with $\|x\|_V=(x^\top Vx)^{1/2}$,
\begin{equation}\label{eq:composite-iss-bound}
\|\xi_t\|_V
\le \rho^t\|\xi_0\|_V+
\sum_{s=0}^{t-1}\rho^{t-1-s}\|d_{s+1}\|_V.
\end{equation}
If $\|d_{t+1}\|_V\le\bar d$, the second term is at most
$\bar d(1-\rho^t)/(1-\rho)$. For a constant fault with $(G-D)f=0$,
zero damping, no model or observation error, and inactive projection, the
origin converges exponentially.
\end{proposition}
\begin{proof}
To obtain \eqref{eq:composite-augmented}, first substitute
$\hat\theta_t=\delta_t+f_t$ into the position equation. For the parameter equation,
subtract $f_{t+1}$ from the update. The coefficient of $f_t$ is
$aC_t+K_tH-I=(a-1)C_t$, and substituting the new position equation supplies
$T_tA$, $-T_tD$ and $T_tb_{t+1}$. Projection adds exactly $p_{t+1}$.
Under \eqref{eq:common-composite-metric},
$\|\mathcal A_t x\|_V\le\rho\|x\|_V$. The triangle inequality followed by
induction proves \eqref{eq:composite-iss-bound}; summing the geometric series
proves the uniform-input bound. In the last stated case every forcing term is
zero.
\end{proof}

Projection cannot simply be dropped from this argument: Euclidean box projection
need not be nonexpansive in a metric with position--parameter cross terms.
One sufficient extension is $V=\operatorname{diag}(V_e,cI_p)$ with $c>0$ and
$f_{t+1}\in\mathcal U$, where $\mathcal U$ is the implemented closed convex box.
Then
$\|\Pi_{\mathcal U}(\bar\theta)-f_{t+1}\|_2
\le\|\bar\theta-f_{t+1}\|_2$, so the same bound holds with the projection term
omitted from $d$, provided the common-metric inequality holds for this block
metric. For a general $V$, a separate bound on $p$ is needed.

For fixed gains, a Schur $\mathcal A$ admits a quadratic contraction metric.
For changing gains, checking each matrix separately is insufficient. For example,
$\left[\begin{smallmatrix}.5&2\\0&.5\end{smallmatrix}\right]$ and
$\left[\begin{smallmatrix}.5&0\\2&.5\end{smallmatrix}\right]$ each have spectral
radius $.5$, but their two-step product has spectral radius $2.25+\sqrt5>1$.
This example is a general switching counterexample, not an assertion that these
two matrices occur in our experiments. The existing candidate screen checks
frozen augmented matrices over the covariance transient and at its limit; it
does not verify \eqref{eq:common-composite-metric}, even for that finite horizon.
Likewise, a fitted $A$ satisfying $A^\top LA\prec L$ need not make the augmented
system contractive. The proposition is conditional on additional, computable
matrix inequalities and valid model/error bounds.

\paragraph{Damping introduces forcing.}
For a constant fault with $(G-D)f=0$, $w=v=p=0$ and fixed gains, positive damping
leaves $d=(0,((a-1)Cf)^\top)^\top$. If $\mathcal A$ is Schur, then
\begin{equation}\label{eq:damped-equilibrium}
\xi_*=(I-\mathcal A)^{-1}
\begin{bmatrix}0\\(a-1)Cf\end{bmatrix}.
\end{equation}
Thus stability generally gives a biased equilibrium rather than exact fault
cancellation. For the scalar prediction-only case $H=1$, $K=k$, $T=0$,
\begin{equation}
\hat\theta_* = \frac{k}{1-a(1-k)}f.
\end{equation}
For $0<k<1$ and $0<a<1$ this attenuates a nonzero fault. With zero damping,
fixed-gain Kalman is the $T=0$ special case; equality of algorithms requires
matching $Q_c,R,P_0$ and projection as well.

\begin{proposition}[Positive tracking gain can destabilize a stable observer]
\label{prop:tracking-limit}
There are an exactly matched scalar plant with contractive healthy dynamics and
valid steady Kalman covariances for which increasing positive tracking gain
destabilizes the implemented local composite recursion.
\end{proposition}
\begin{proof}
Take $A=.8$, $G=D=H=L=1$, $\lambda=0$, $\Delta t=1$, $R=1$, $Q_c=.05$ and
$P_t=P_t^+=.2$. Then $P_t^-=.25$, $K_t=.2$, and $T_t=.2\eta$, all exactly
consistent with the covariance recursion. For inactive projection,
\begin{equation}
\mathcal A(\eta)=
\begin{bmatrix}.8&-1\\.16\eta&.8-.2\eta\end{bmatrix},\qquad
\operatorname{tr}\mathcal A=1.6-.2\eta,\quad
\det\mathcal A=.64.
\end{equation}
The second-order Schur conditions
$1-\det\mathcal A>0$, $1-\operatorname{tr}\mathcal A+\det\mathcal A>0$ and
$1+\operatorname{tr}\mathcal A+\det\mathcal A>0$ hold, for $\eta\ge0$, precisely
when $\eta<16.2$. Kalman alone ($\eta=0$) is stable, whereas $\eta=20$ gives
eigenvalues $-1.2\pm\sqrt{.8}$, one of magnitude greater than one. The same
local instability holds around the zero-fault equilibrium with any positive
box bound, since projection is inactive in a neighborhood of that equilibrium.
\end{proof}

The threshold is for this scalar example's rate $\eta$ and unit time step; it
is not a threshold for the ALOHA rate or its dimensionless active-strength setting.
This counterexample concerns the algorithm class; it is not a diagnosis of the
screened experimental configurations. Even within a stable range, the extra
term can introduce sensitivity to reference-model error or redistribute gains
across coordinates. Neither a positive gradient direction nor an empirical
contraction check implies a universal advantage over Kalman or the other
estimators.

\subsection{What would permit a task-success conclusion?}

\begin{proposition}[Transfer through an assumed robust trajectory tube]
\label{prop:task-tube}
Fix a successful nominal complete execution $x^0_{0:N}$. Suppose every execution
with $\max_{t\le N}\|x_t-x_t^0\|\le\epsilon_{\rm task}$ also succeeds.
Suppose, independently, the full-execution discrepancy satisfies
$\|x_t-x_t^0\|\le c\|\xi_t\|_V+\epsilon_{\rm lift}$ for all $t\le N$, and
the assumptions of Proposition~\ref{prop:composite-iss} hold with
$\|d_t\|_V\le\bar d$. Then the corrected execution succeeds whenever
\begin{equation}
c\max\!\left\{\|\xi_0\|_V,\frac{\bar d}{1-\rho}\right\}
+\epsilon_{\rm lift}\le\epsilon_{\rm task}.
\end{equation}
\end{proposition}
\begin{proof}
The bound in \eqref{eq:composite-iss-bound} is at most a convex combination of
$\|\xi_0\|_V$ and $\bar d/(1-\rho)$, hence no larger than their maximum.
The assumed lifting inequality places the complete execution in the successful
tube.
\end{proof}

Both additional assumptions are substantive. The complete execution must include
the variables needed to determine success; a six-position error does not bound
unobserved velocities, contact states, other joints, or changes in the policy's
observations. The fitted reference driven by each rollout's own raw actions is
also not the common successful nominal execution assumed here. Neither the
robust tube nor the lifting bound was established by our experiments. This
conditional result therefore states what is missing for task-level transfer,
not a success guarantee or a claim that any estimator must rank first.
